\chapter{Симметрийный запрет аффинно-Hodge-селектора пространства копий}
\label{ch:v7-affine-hodge-copy-selector-no-go-gate}

\section{Точный вопрос}

Четырёхвершинный носитель оставляет в каждом лептонном секторе трёхмерное
ядро прямоугольной семейной массы
\begin{equation}
 M:\mathbb C^6\longrightarrow\mathbb C^3,
 \qquad \dim\ker M=3.
 \label{eq:v7-copy-selector-rectangular-mass}
\end{equation}
Проверим, способен ли уже принятый аффинно-Hodge-родитель выбрать положение
этого ядра среди старой и новой копий без второго оператора.

\section{Грассманиан ядер}

Для полного ранга проектор ядра равен
\begin{equation}
 P_{\ker M}=I_6-M^\dagger(MM^\dagger)^{-1}M,
 \qquad \rank P_{\ker M}=3.
 \label{eq:v7-copy-selector-kernel-projector}
\end{equation}
Поэтому возможные лёгкие пространства образуют
\begin{equation}
 P_{\ker M}\in\operatorname{Gr}_{\mathbb C}(3,6),
 \qquad
 \dim_{\mathbb R}\operatorname{Gr}_{\mathbb C}(3,6)=18.
 \label{eq:v7-copy-selector-grassmannian}
\end{equation}
Описание прямоугольных полу-унитарных матриц через многообразие Штифеля и
грассманиан стандартно \cite{v7copy-grassmann}. В проекте оно означает, что
правильная размерность ядра ещё не выбирает физические смеси.

\section{Точная копийная орбита}

Разложим шестимерный вход как
$\mathbb C^3_{\rm fam}\otimes\mathbb C^2_{\rm copy}$. Для любого
$U\in U(2)_{\rm copy}$ положим
\begin{equation}
 M\longmapsto M^U=M(U\otimes I_3).
 \label{eq:v7-copy-selector-orbit}
\end{equation}
Тогда
\begin{equation}
 M^U(M^U)^\dagger=MM^\dagger,
 \qquad
 P_{\ker M^U}=(U\otimes I_3)^\dagger
 P_{\ker M}(U\otimes I_3).
 \label{eq:v7-copy-selector-invariant-gram-moving-kernel}
\end{equation}
Следовательно, сингулярные числа и любой Hodge-функционал от
$MM^\dagger$ постоянны, тогда как само лёгкое подпространство движется.

Явный аудит использует $M_0=(I_3\;0)$ и семнадцать вращений от нуля до
$\pi/2$. Он получает
\begin{equation}
 \Delta\mathcal S_{\rm Hodge}<10^{-24},
 \qquad
 \max\|P_{\ker M^U}-P_{\ker M_0}\|_F=\sqrt6,
 \label{eq:v7-copy-selector-orbit-audit}
\end{equation}
а перекрытие ядра со старой копией проходит весь интервал $[0,3]$.
Гессиан действия вдоль этой орбиты равен нулю.

\section{Почему аффинный проектор не помогает}

Аффинный подъём имеет вид
\begin{equation}
 P_{\rm aff}=P_3\otimes I_2,
 \qquad \rank P_3=3.
 \label{eq:v7-copy-selector-affine-lift}
\end{equation}
Поэтому для всех копийных вращений
\begin{equation}
 [P_3\otimes I_2,I_4\otimes U]=0.
 \label{eq:v7-copy-selector-affine-commutator}
\end{equation}
$P_3$ выбирает стандартный семейный триплет внутри аффинной четвёрки, но
действует единично на метке старой/новой копии.

То же препятствие видно из коммутанта. На совпадающих представлениях
алгебра действует как
\begin{equation}
 \pi(a)=\pi_0(a)\otimes I_2,
 \qquad
 I\otimes M_2(\mathbb C)\subset\pi(\mathcal A_{\rm SM})'.
 \label{eq:v7-copy-selector-copy-commutant}
\end{equation}
Машинный аудит проверяет четыре независимые матричные единицы этого
коммутанта с нулевым остатком. Хиральная градуировка, Real-сопряжение и
класс $15$ также не действуют нецентрально внутри пары одинаковой
хиральности и калибровочного типа.

\section{Литературная граница}

Квадрат нормы отображения момента на пространстве представлений колчана
может стратифицировать орбиты, но устойчивость требует отдельного характера
или параметра \cite{v7copy-quiver}. В текущей архитектуре такого
copy-space параметра нет. Его подстановка как
$\operatorname{diag}(1,-1)$ означала бы заранее выбрать старую и новую
копии, то есть записать ответ в действие.

\begin{theorem}[Аффинно-Hodge no-go пространства копий]
Тривиальный подъём существующего аффинно-Hodge-родителя на
$\mathbb C^2_{\rm copy}$ инвариантен относительно полного $U(2)$ копий.
Он фиксирует ранг массы и размерность лёгкого ядра, но не способен выбрать
его ориентацию, различить старые и новые совпадающие бимодули или выделить
шесть целевых рёбер из одиннадцати разрешённых.
\end{theorem}

\begin{nogocriterion}[Запрет ручной копийной градуировки]
Нельзя добавлять матрицу $\operatorname{diag}(1,-1)$, проектор старой копии
или пять нулевых юкавских блоков как входные данные. Следующий кандидат
должен вывести второй нецентральный оператор из ориентированной алгебры
путей, Morita-угла либо другого независимо существующего носителя.
\end{nogocriterion}

\begin{protocol}
\begin{enumerate}
 \item размерность ядра: \textbf{выбирается};
 \item ориентация ядра: \textbf{не выбирается};
 \item Hodge-действие вдоль $U(2)$-орбиты: \textbf{постоянно};
 \item проектор ядра вдоль орбиты: \textbf{движется};
 \item аффинный $P_3$ нарушает $U(2)$ копий: \textbf{нет};
 \item коммутант содержит $M_2(\mathbb C)$: \textbf{да};
 \item шесть из одиннадцати рёбер выбраны: \textbf{нет};
 \item старый родитель как copy-selector: \textbf{закрыт};
 \item следующий гейт: \textbf{происхождение нецентрального оператора из
 алгебры путей или Morita-угла}.
\end{enumerate}
\end{protocol}

Машинный сертификат сохранён в
\path{s2t_v7_affine_hodge_copy_selector_no_go_gate_results.json}.

\begin{thebibliography}{9}
\bibitem{v7copy-grassmann}
T.~Bendokat, R.~Zimmermann, P.-A.~Absil,
\emph{A Grassmann Manifold Handbook: Basic Geometry and Computational Aspects},
Advances in Computational Mathematics 50 (2024), article 6;
arXiv:2011.13699.
\bibitem{v7copy-quiver}
M.~Harada, G.~Wilkin,
\emph{Morse theory of the moment map for representations of quivers},
Geometriae Dedicata 150 (2011), 307--353; arXiv:0807.4734.
\end{thebibliography}